% ============================================================================
% Part 3 --- The Rust reference verifier
% ============================================================================
\chapter{The Rust Reference Verifier}
\label{ch:rust}

This chapter walks the Rust verifier that the code generator must reproduce.
It is the \emph{source of truth}: if the generated Solidity disagrees with
anything here, the Solidity is wrong. All line numbers refer to the checkout
at the time of writing.

\section{Module map}
\label{sec:rust-modules}

\begin{longtable}{@{}p{0.34\linewidth}p{0.60\linewidth}@{}}
\toprule
\textbf{Module} & \textbf{Role} \\
\midrule
\endfirsthead
\toprule
\textbf{Module} & \textbf{Role} \\
\midrule
\endhead
\rust{plonk/verifier.rs} (696 L) & Top-level: \rust{parse\_\allowbreak{}trace},
  \rust{verify\_\allowbreak{}algebraic\_\allowbreak{}constraints}, \rust{prepare}. Drives the whole
  transcript schedule and assembles the PCS query list. \\
\rust{plonk/mod.rs} & \rust{partially\_\allowbreak{}evaluate\_\allowbreak{}identities} (L488):
  builds the ordered identity-evaluation list; \rust{hash\_\allowbreak{}into} (L283):
  VK digest absorption. \\
\rust{plonk/\allowbreak{}linearization/\allowbreak{}verifier.rs} (104 L) &
  \rust{compute\_\allowbreak{}linearization\_\allowbreak{}commitment}: Eq.~\eqref{eq:linearization}. \\
\rust{plonk/\allowbreak{}permutation.rs} & \rust{expressions} (L187): the four
  permutation identity families \eqref{eq:perm1}--\eqref{eq:perm-right}. \\
\rust{plonk/\allowbreak{}permutation/\allowbreak{}verifier.rs} (198 L) & Reads $z$ commitments,
  evaluates $z$/$z_{\text{next}}$/$z_{\text{last}}$, emits PCS queries. \\
\rust{plonk/logup.rs} & \rust{expressions} (L400): helper \eqref{eq:lookup-helper}
  and accumulator \eqref{eq:lookup-acc} identities. \\
\rust{plonk/\allowbreak{}logup/\allowbreak{}verifier.rs} (212 L) & Reads multiplicities, helpers,
  accumulator; evaluates and emits PCS queries. \\
\rust{plonk/trash.rs} & \rust{expressions} (L58): trash identity
  \eqref{eq:trash}. \\
\rust{plonk/\allowbreak{}trash/\allowbreak{}verifier.rs} (75 L) & Reads trash commitments, evaluates,
  emits PCS queries. \\
\rust{transcript/mod.rs} (184 L) & \rust{Transcript} trait:
  \rust{common}/\rust{read}/\rust{squeeze\_\allowbreak{}challenge}/\rust{assert\_\allowbreak{}empty}. \\
\rust{transcript/\allowbreak{}implementors.rs} (341 L) & Keccak256 hash;
  \rust{sample\_\allowbreak{}keccak\_\allowbreak{}digest\_\allowbreak{}be\_\allowbreak{}mod\_\allowbreak{}r} (L19); G1/Fq \rust{to\_\allowbreak{}input}. \\
\rust{poly/kzg/mod.rs} (738 L) & \rust{multi\_\allowbreak{}prepare} (L321): the KZG
  multi-open reduction (Sec.~\ref{sec:pcs-math}). \\
\rust{poly/kzg/utils.rs} (187 L) & \rust{construct\_\allowbreak{}intermediate\_\allowbreak{}sets}
  (L37), \rust{compute\_\allowbreak{}dummy\_\allowbreak{}queries} (L150). \\
\rust{poly/kzg/msm.rs} & \rust{DualMSM}, final pairing check
  \eqref{eq:final-pairing-dual}. \\
\rust{poly/domain.rs} & \rust{rotate\_\allowbreak{}omega} (L410), \rust{l\_\allowbreak{}i\_\allowbreak{}range}
  (L448), \rust{get\_\allowbreak{}quotient\_\allowbreak{}poly\_\allowbreak{}degree} (L476). \\
\rust{utils/\allowbreak{}arithmetic.rs} & \rust{truncate} (L247),
  \rust{truncated\_\allowbreak{}powers} (L255), \rust{powers}, \rust{inner\_\allowbreak{}product},
  \rust{msm\_\allowbreak{}inner\_\allowbreak{}product}, \rust{lagrange\_\allowbreak{}interpolate} (L172),
  \rust{eval\_\allowbreak{}polynomial} (L53). \\
\rust{poly/query.rs} & \rust{VerifierQuery}, \rust{CommitmentReference},
  \rust{CommitmentLabel}. \\
\bottomrule
\end{longtable}

\section{\rust{parse\_\allowbreak{}trace} --- reading the proof into a trace}
\label{sec:parse-trace}

\rust{parse\_\allowbreak{}trace} (L24) consumes the proof byte stream and produces a
\rust{VerifierTrace} struct holding every commitment and challenge in
transcript order. This is the function whose schedule Table~\ref{tab:schedule}
describes. Its shape (excerpt, L24--120):

\begin{lstlisting}[style=rust]
pub fn parse_trace<F, CS: PolynomialCommitmentScheme<F>, T: Transcript>(
    vk: &VerifyingKey<F, CS>,
    instances: &[&[&[F]]],
    transcript: &mut T,
) -> Result<VerifierTrace<F, CS>, Error> {
    vk.hash_into(transcript);                       // (1) vk_digest absorb
    // committed-instance identity absorb, num_instances + instance absorbs
    // per-phase advice reads + challenge squeezes
    // theta, multiplicities, beta, gamma, permutation product comms,
    // lookup helpers + accumulator, trash_challenge, trash comms, y,
    // quotient limbs, x, ...
}
\end{lstlisting}

The trace struct (\rust{plonk/traces.rs}) carries:
\begin{lstlisting}[style=rust]
pub struct VerifierTrace<F, CS> {
    pub advice_commitments: Vec<Vec<CS::Commitment>>, // per user phase
    pub lookups: Vec<LookupTrace<CS>>,               // m, helpers, acc
    pub trashcans: Vec<CS::Commitment>,
    pub permutations: Vec<CS::Commitment>,           // z products
    pub challenges: Vec<F>,                          // user challenges
    pub beta: F, pub gamma: F, pub theta: F,
    pub trash_challenge: F, pub y: F,
}
\end{lstlisting}

\section{\rust{verify\_\allowbreak{}algebraic\_\allowbreak{}constraints} --- the main body}
\label{sec:verify-alg}

This is the function the generated Solidity mirrors most directly. Its
signature (L238) takes the VK, the parsed trace, the committed instances,
the public instances, and the transcript, and returns a PCS verification
guard. The body proceeds in the following stages.

\subsection{Instance validation}
\begin{lstlisting}[style=rust]
if committed_instances.is_empty() {
    return Err(Error::InvalidInstances);
}
let nb_committed_instances = committed_instances[0].len();
let num_proofs = instances.len();
\end{lstlisting}
The supported profile additionally requires exactly one committed and one
non-committed instance column; the codegen enforces this at build time
(Section~\ref{sec:profile}) so the generated contract never has to branch
on it.

\subsection{Quotient limbs and the evaluation challenge}
\begin{lstlisting}[style=rust]
let nb_quotient_coms = vk.domain.get_quotient_poly_degree();
let quotient_limb_coms = read_n(transcript, nb_quotient_coms)?;
let x: F = transcript.squeeze_challenge();
let splitting_factor = x.pow_vartime([vk.n() - 1]);
let xn = splitting_factor * x;
\end{lstlisting}
This is rows 15--16 of Table~\ref{tab:schedule} and the splitting factor
of Eq.~\eqref{eq:linearization}. \texttt{single-h-commitment} collapses
the limbs to one commitment (the \texttt{outer-single-h-commitment}
profile).

\subsection{Instance evaluation}
\begin{lstlisting}[style=rust]
let l_i_s = &vk.domain.l_i_range(
    x, xn,
    -max_rotation..max_instance_len as i32 + min_rotation.abs(),
);
\end{lstlisting}
Implements \eqref{eq:lagrange}--\eqref{eq:instance-eval}. The
\texttt{l\_\allowbreak{}i\_\allowbreak{}range} helper uses the barycentric weight $1/n$
(\rust{domain.rs} L448).

\subsection{Partial identity evaluation}
\begin{lstlisting}[style=rust]
let expressions = partially_evaluate_identities(
    vk, &fixed_evals, &instance_evals, &advice_evals,
    permutations_evaluated.iter().map(|e| &e.sets),
    lookups_evaluated.iter().map(|e| e.iter().map(|inner| &inner.evaluated)),
    trashcans_evaluated.iter().map(|e| e.iter().map(|inner| &inner.evaluated)),
    &permutations_common, x, xn, beta, gamma, theta, trash_challenge,
    &challenges,
);
\end{lstlisting}
Returns \rust{Vec<(Option<usize>, F)>}: each identity's evaluation tagged
with its simple-selector column index (\texttt{None} = fully evaluated).
This is the ordered list $\bigl(e_i(x)\bigr)$ of \eqref{eq:nu-y}.

\subsection{Linearization and query assembly}
\begin{lstlisting}[style=rust]
let lin_com = compute_linearization_commitment(
    expressions, vk, x, &y, &xn, &splitting_factor, &quotient_limb_coms);
// ... queries = advice ++ committed-instance ++ perm-z ++ lookup ++ trash
//     ++ fixed(non-simple) ++ perm-common ++ lin_com
CS::multi_prepare(&queries, transcript).map_err(|_| Error::Opening)
\end{lstlisting}
The query list order (L577--652) is exactly the PCS query order the
codegen's \cg{ProtocolPlan} reproduces (Section~\ref{sec:protocol-plan}).

\section{\rust{partially\_\allowbreak{}evaluate\_\allowbreak{}identities}}
\label{sec:partial-eval}

In \rust{plonk/mod.rs} (L488), this function:
\begin{enumerate}[leftmargin=1.4em,itemsep=1pt]
  \item computes $\ell_{\mathsf{last}}, \ell_{\mathsf{blind}}, \ell_0$
        from \eqref{eq:llast}--\eqref{eq:l0};
  \item for each gate: evaluates the gate polynomial with the fixed/advice/
        instance evals, and tags it with the simple-selector index via
        \rust{gate.queried\_\allowbreak{}selectors().filter(is\_\allowbreak{}simple).next()};
  \item appends permutation expressions (tagged \texttt{None}),
        lookup expressions (\texttt{None}), trash expressions (\texttt{None}).
\end{enumerate}
The ordering gates $\to$ permutation $\to$ lookup $\to$ trash fixes the
$y$-weights in \eqref{eq:nu-y}.

\section{\rust{compute\_\allowbreak{}linearization\_\allowbreak{}commitment}}
\label{sec:lin-com}

The full function was quoted in Section~\ref{sec:linearization}; the key
excerpt (L63--94):
\begin{lstlisting}[style=rust]
identities_points.extend(quotient_limb_commitments);
let mut splitting_pow = F::ONE - *xn;
for _ in 0..quotient_limb_commitments.len() {
    identities_scalars.push(splitting_pow);
    splitting_pow *= splitting_factor;
}
let mut grouped_points: BTreeMap<Option<usize>, F> = BTreeMap::new();
let mut y_pow = F::ONE;
expressions.iter().rev().for_each(|(col_idx, eval)| {
    *grouped_points.entry(*col_idx).or_insert(F::ZERO) += y_pow * eval;
    y_pow *= y;
});
let mut expected_eval = F::ZERO;
grouped_points.into_iter().for_each(|(col_idx, eval)| {
    match col_idx {
        Some(col_idx) => { identities_points.push(&vk.fixed_commitments[col_idx]);
                          identities_scalars.push(eval); }
        None => { expected_eval -= eval; }
    }
});
\end{lstlisting}
This is literally \eqref{eq:linearization} and \eqref{eq:buckets}: the
quotient limbs get scalars $(1-x^n)\cdot x_{\mathsf{split}}^i$; the
reverse iteration with running $y^{\text{pow}}$ is the Horner form of
\eqref{eq:nu-y}; \texttt{None} entries negate into the expected eval
(Definition~\ref{def:neg-numer}).

\section{Permutation identities}
\label{sec:rust-perm}

\rust{permutation.rs::expressions} (L187) builds the four families
\eqref{eq:perm1}--\eqref{eq:perm-right}. The chunk running constant is
\begin{lstlisting}[style=rust]
let current_delta = beta * x * delta.pow_vartime([(chunk_index * chunk_len) as u64]);
\end{lstlisting}
with \rust{chunk\_\allowbreak{}len = cs\_\allowbreak{}degree - 2}. The verifier side
(\rust{permutation/\allowbreak{}verifier.rs}) reads \rust{read\_\allowbreak{}product\_\allowbreak{}commitments}
(one $z$ per set), evaluates $z_{\text{eval}}, z_{\text{next\_\allowbreak{}eval}}$ and
$z_{\text{last\_\allowbreak{}eval}}$ (the latter only when more than one set exists),
and emits PCS queries at $x$, $x_{\mathsf{next}} = \omega x$,
$x_{\mathsf{last}} = \omega^{-(b+1)}x$.

\section{Lookup (LogUp) identities}
\label{sec:rust-lookup}

\rust{logup.rs::expressions} (L400): compression
$\mathsf{fold}(\mathsf{acc}\cdot\theta + \mathsf{eval})$; boundary
$(\ell_0 + \ell_{\mathsf{last}})\cdot Z_{\mathsf{eval}}$; helper
constraint \eqref{eq:lookup-helper} with partial products via
$\bigl(\prod_k (f_k+\beta)\bigr)\cdot(f_j+\beta)^{-1}$; accumulator
constraint \eqref{eq:lookup-acc} times the active-row mask
$1-(\ell_{\mathsf{last}}+\ell_{\mathsf{blind}})$. The verifier
(\rust{logup/verifier.rs}) reads multiplicities, then per lookup
$\mathsf{nb\_\allowbreak{}chunks}$ helper commitments $+\,1$ accumulator, evaluates
$m$, helpers, $z$, $z_{\text{next}}$, and emits PCS queries.

\section{Trash identities}
\label{sec:rust-trash}

\rust{trash.rs::expressions} (L58): compressed $=$ fold of constraint
evals with $\mathsf{acc}\cdot\mathsf{trash\_\allowbreak{}challenge} + \mathsf{eval}$;
identity $=$ compressed $-\,(1-q)\cdot\mathsf{trash\_\allowbreak{}eval}$
\eqref{eq:trash}. The verifier reads committed trash, evaluates, emits
queries.

\section{Transcript implementation}
\label{sec:rust-transcript}

\rust{transcript/\allowbreak{}implementors.rs} defines the Keccak256 hash (L186):
\begin{lstlisting}[style=rust]
fn squeeze(&mut self) -> [u8; 32] {
    let digest = self.state.clone().finalize();   // clone, don't consume
    self.state.reset();
    self.state.update(digest.as_bytes());        // reseed with digest
    *digest.as_ref()
}
\end{lstlisting}
with sampling \rust{sample\_\allowbreak{}keccak\_\allowbreak{}digest\_\allowbreak{}be\_\allowbreak{}mod\_\allowbreak{}r} (L19) interpreting
the digest big-endian and reducing mod $r$. G1 \rust{to\_\allowbreak{}input} (L283)
emits the 128-byte EIP-2537 uncompressed form \eqref{eq:g1-encoding};
Fq \rust{to\_\allowbreak{}input} (L315) emits the reversed (big-endian) repr. These
are the exact byte-level contracts the generated parser must match.

\subsection{Production G1 \texttt{to\_input} (Keccak256)}
\label{sec:rust-to-input}

The verifier's G1 absorption uses
\rust{Hashable<Keccak256> for midnight\_\allowbreak{}curves::G1Projective}
(\texttt{implementors.rs} L264--297):
\begin{lstlisting}[style=rust]
fn to_input(&self) -> Vec<u8> {
    let aff = G1Affine::from(self);
    let mut out = vec![0u8; 128];
    if !bool::from(aff.is_identity()) {
        let raw = <G1Affine as UncompressedEncoding>::to_uncompressed(&aff);
        let bytes: &[u8] = raw.as_ref();
        out[16..64].copy_from_slice(&bytes[0..48]);    // x_be
        out[80..128].copy_from_slice(&bytes[48..96]);  // y_be
    }
    out   // identity => 128 zero bytes
}
\end{lstlisting}
This is the precise contract the generated
\sol{common\_\allowbreak{}uncompressed\_\allowbreak{}g1} mirrors: the identity
encodes as 128 zero bytes (EIP-2537 convention, \emph{not} the BLS-spec
\texttt{0x40}-leading form), and non-identity points place the 48-byte
big-endian coordinates at offsets 16 and 80. The wire format
(\rust{to\_\allowbreak{}bytes}) is unchanged --- 48-byte compressed --- so the
uncompressed form is a transcript-only encoding.

\section{Top-level entry: \rust{prepare}}
\label{sec:rust-prepare}

The public verifier entry point
\rust{plonk/\allowbreak{}verifier.rs::prepare} (L659) is a thin composition of the
two functions already covered:
\begin{lstlisting}[style=rust]
pub fn prepare<F, CS, T>(vk, committed_instances, instances, transcript)
    -> Result<CS::VerificationGuard, Error>
{
    let trace = parse_trace(vk, committed_instances, instances, transcript)?;
    verify_algebraic_constraints(vk, trace, committed_instances, instances, transcript)
}
\end{lstlisting}
It returns the PCS \rust{VerificationGuard} (the \rust{DualMSM}) rather
than a boolean; the caller invokes \rust{Guard::verify} to perform the
final pairing. The generated Solidity collapses this three-step flow
(\rust{parse\_trace} $\to$ \rust{verify\_\allowbreak{}algebraic\_\allowbreak{}constraints} $\to$
\rust{DualMSM::check}) into one \texttt{verifyProof} that returns
\texttt{true} or reverts --- the guard is never materialized as a value.

\section{The \rust{DualMSM} pairing accumulator}
\label{sec:rust-dualmsm}

The PCS guard is \rust{DualMSM<E>} (\texttt{poly/kzg/msm.rs}), a
two-channel MSM accumulator. Its \rust{check} (L284) is the final
pairing:
\begin{lstlisting}[style=rust]
pub fn check(self, params: &ParamsVerifierKZG<E>) -> bool {
    let left = if self.left.scalars.len() == 1 && self.left.scalars[0] == F::ONE {
        self.left.bases[0]            // single scalar-1 term: skip MSM
    } else {
        self.left.eval()
    };
    let right = self.right.eval();
    let terms = &[(&left.into(), &params.s_g2_prepared),
                 (&right.into(), &params.n_g2_prepared)];
    bool::from(E::multi_miller_loop(&terms[..])
        .final_exponentiation().is_identity())
}
\end{lstlisting}
This is \eqref{eq:final-pairing-dual}:
$e(\mathsf{left},[s]_2)\cdot e(\mathsf{right},-[1]_2)=1$ (with
\texttt{n\_g2\_prepared} $=-[1]_2$). The \texttt{Guard::verify} trait
impl (L228) wraps it as
\texttt{self.check(params).then\_\allowbreak{}some(()).ok\_\allowbreak{}or(Error::OpeningError)}.

\begin{itemize}
  \item \textcolor{mathpurple}{\textbf{[NOTE] Single-scalar-1 fast path.}}
        When the left channel has exactly one term with scalar $1$, the
        MSM is skipped and the base is used directly. The generated
        \sol{ec\_pairing} does not replicate this micro-optimization
        (it always runs the fused MSM); the result is identical, only
        the Rust saves one MSM. No faithfulness impact.
  \item \textcolor{tracegreen}{\textbf{[OK] Channel semantics.}}
        \texttt{left} carries the $\pi$-side (paired with $[s]_2$) and
        \texttt{right} the combined side (paired with $-[1]_2$),
        matching the \texttt{PAIRING\_LHS}/\texttt{PAIRING\_RHS}
        mapping documented in \S\ref{sec:ec-pairing}.
\end{itemize}

\section{KZG \rust{multi\_\allowbreak{}prepare}}
\label{sec:rust-multiprep}

Quoted in full in Section~\ref{sec:pcs-math}; the load-bearing excerpt is
the $f_{\mathsf{eval}}$ fold (L471--482):
\begin{lstlisting}[style=rust]
let f_eval = point_sets.iter().zip(q_eval_sets.iter())
    .zip(q_evals_on_x3.iter()).rev().fold(E::Fr::ZERO,
    |acc_eval, ((points, evals), proof_eval)| {
        let r_poly = lagrange_interpolate(points, evals);
        let r_eval = eval_polynomial(&r_poly, x3);
        let den = points.iter().fold(E::Fr::ONE, |acc, p| acc * &(x3 - p));
        let eval = (*proof_eval - &r_eval) * den.invert().unwrap();
        acc_eval * &(x2) + &eval
    });
\end{lstlisting}
which is \eqref{eq:feval}. The final commitment and scalar (L488--527)
are \eqref{eq:finalcom}--\eqref{eq:v}; the dual-MSM assembly (L540--558)
is \eqref{eq:final-pairing-dual}.

\section{Arithmetic helpers}
\label{sec:rust-arith}

\begin{itemize}[leftmargin=1.4em,itemsep=1pt]
  \item \rust{truncate} (L247): keeps the lower 16 little-endian bytes
        ($=128$ bits).
  \item \rust{truncated\_\allowbreak{}powers} (L255): $\mathsf{truncate}(x^i)$ with a
        full-precision running accumulator.
  \item \rust{lagrange\_\allowbreak{}interpolate} (L172): builds $r_s$ of
        \eqref{eq:feval}.
  \item \rust{eval\_\allowbreak{}polynomial} (L53): Horner evaluation.
  \item \rust{inner\_\allowbreak{}product}/\rust{msm\_\allowbreak{}inner\_\allowbreak{}product} (L263/L275):
        Eqs.~\eqref{eq:qcom}--\eqref{eq:v}.
\end{itemize}

\section{Trace instrumentation}
\label{sec:rust-trace}

Under the \texttt{solidity-verifier-trace} feature, every semantic value is
recorded with a stable trace id (\rust{plonk/solidity\_\allowbreak{}trace.rs}):
\boxid{1 vk\_\allowbreak{}digest}, \boxid{2 num\_\allowbreak{}instances}, \boxid{3 k},
\boxid{4 n\_\allowbreak{}inv}, \boxid{5 omega}, \boxid{6 omega\_\allowbreak{}inv}, \boxid{7 theta},
\boxid{8 beta}, \boxid{9 gamma}, \boxid{10 y}, \boxid{11 x},
\boxid{12 trash\_\allowbreak{}challenge}, \boxid{13 x1}, \boxid{14 x2}, \boxid{15 x3},
\boxid{16 x4}, \boxid{17 x\_\allowbreak{}n}, \boxid{18 x\_\allowbreak{}n\_\allowbreak{}minus\_\allowbreak{}1\_\allowbreak{}inv},
\boxid{19 l\_\allowbreak{}last}, \boxid{20 l\_\allowbreak{}blind}, \boxid{21 l\_\allowbreak{}0},
\boxid{22 instance\_\allowbreak{}eval}, \boxid{23 quotient\_\allowbreak{}eval}, \boxid{24 quotient},
\boxid{25 f\_\allowbreak{}com}, \boxid{26 pi}, \boxid{29/30 acc lhs/rhs},
\boxid{31 f\_\allowbreak{}eval}, \boxid{32 v}, \boxid{33 final\_\allowbreak{}com},
\boxid{34 linearization com}, \boxid{35 final result}, \boxid{36 nu\_\allowbreak{}y(x)};
bases: \boxid{1000} user challenges, \boxid{10000} proof commitments,
\boxid{20000} proof evals, \boxid{30000} quotient identities,
\boxid{40000} q\_\allowbreak{}com, \boxid{41000} point sets, \boxid{60000} selector
folds. These ids are the backbone of the differential equivalence test
(Chapter~\ref{ch:faithfulness}).
